\section{Authoring and Traveler Workflows}
\label{sec:features}

\subsection{Creator workspace and trip lifecycle}
The Creator opens a local trip library. An agent selects a trip or creates one with a title, date range, description, and banner. The stable design adds destination tags, a primary IANA time zone, currency, language, and a local traveler directory. An agent can keep contact, dietary, emergency, and passport-expiry information in that directory. The Creator compares the passport expiry date with the trip end date and flags validity under six months for the agent to review. The trip view exposes day schedules and a separate Information and Documents area. The agent can find exported copies in a library of timestamped files. The editable workspace remains distinct from those exports.

Trip dates supply the initial sequence of day records. Each day belongs to the trip and has its own date, title, and ordered events. In the evaluated Creator, extending the date range adds days. Shortening it removes empty days beyond the new range and retains days containing events. The stable Creator places the trip and day navigator beside the timeline and uses an inspector for the selected event or route. Moving a card across days changes its day association. The scheduler sorts timed cards by start time and uses display order for ties. It converts each timestamp to an instant using its recorded IANA zone. Two timed events overlap when each begins before the other ends. The agent sees a conflict and can adjust the relevant card in the inspector. Untimed notes retain a display position. Flight departure and arrival each carry their own time zone. The beta allows day reordering, but its export mapper defaults event time zones to UTC and can supply default times for untimed records.

\subsection{Typed event cards}
An event card gives an agent an editing surface and a traveler a summary of what to do. The stable data model assigns each event an identifier, a day, a display order, and a type. The renderer selects the fields appropriate to the type and resolves optional attachments through the asset index. The six cases below describe the intended fields and actions. The beta has typed card groups and a partial mapping into these records.

\begin{itemize}
\item \emph{Flight}. The agent enters the airline, flight number, airports, local departure and arrival times, terminals, seats, baggage, and reference number. The traveler sees both ends of the journey and can open the linked boarding pass or copy the booking reference.
\item \emph{Lodging}. The agent records check-in and check-out, property contacts, room and meal details, and confirmation. The traveler can see the property on a map, call it, or open its voucher and photos.
\item \emph{Transportation}. The agent chooses a mode and enters carrier, pickup and drop-off, coach or seat, and instructions. The traveler sees the pickup point and time next to the ticket. The beta retains display subcategories, while part of its structured transport mapping defaults to driving.
\item \emph{Activity}. The card holds the meeting point, start time, duration, operator, and instructions. The traveler can open the location or ticket from the card.
\item \emph{Dining}. The card holds the restaurant, reservation time, cuisine or dietary notes, and any menu or reservation file. Its location opens through the map action.
\item \emph{Information and advisory}. A note carries formatted instructions with an optional category, severity, checklist, or supporting file. Examples include visa reminders, packing instructions, and emergency contacts.
\end{itemize}

\subsection{Location cards and navigation}
The Creator can attach more than one location card to an event. A transfer can have separate pickup and drop-off cards. The parser extracts a name and coordinates where possible. Short links may require a redirect lookup. Missing coordinates or place names may trigger a Nominatim request. Each saved card retains the available name, coordinate pair, and map URL.

The beta Viewer opens a saved map link or creates a place search from a coordinate pair or name. For a day route, it traverses events in display order and collects one query per usable location card. It removes consecutive duplicate queries. A single stop opens a search. Two or more stops become the origin, destination, and intervening waypoints in a Google Maps directions URL \cite{maps}. This workflow delegates route computation to the map provider.

The stable Creator walks the day's ordered stops to build route connectors between consecutive locations. The agent can request an online ground route or supply manual or geodesic geometry for a leg. Each connector records its endpoints, mode, encoded polyline, distance, and duration. The agent can also draw an area or add a point of interest with opening hours and notes. The Viewer filters connectors, numbered stops, attraction pins, and polygons by the selected day. MapLibre draws these overlays above a basemap and an optional PMTiles street slice. Selecting a route focuses it on the map. Selecting an attraction opens its details. Embedded tiles and saved geometry support local display during the trip.

\subsection{Images, documents, and information}
The agent can add a trip banner, place event photos through the file picker or clipboard, and attach documents to either an event or an entry in Information and Documents. The two document entry points support different tasks. A rail ticket belongs beside its departure event. General instructions and emergency contacts belong at trip level. The editor stores description formatting with the record, so the Viewer can reconstruct headings, lists, and inline emphasis from serialized text. Rendering parity for every desktop control has yet to be evaluated.

The beta workspace and exporter can deduplicate identical documents and photos by hashing their bytes. Records can retain separate titles and filenames while sharing a stored blob. The exporter embeds readable local assets and records their digest, offset, and byte length. A banner can travel in its own field. Some external image URLs remain links, and a document whose source file is missing can retain metadata without usable bytes. The exporter also prepends synthesized note events for certain trip-level document records. The stable design keeps trip-level material in its own collection and renders it independently of the scheduled cards.

The Viewer resolves an embedded image by its asset reference and displays it in a swipeable, zoomable gallery. A document card resolves the corresponding bytes and opens the selected file. The beta can hand a temporary local copy to an installed Android handler. The stable Creator resizes eligible photos to a 1,920-pixel maximum edge and encodes them as WebP or JPEG before hashing the packaged bytes. Ticket and barcode legibility constrains that conversion. The Viewer looks up the asset by digest and sends the bytes to the offline gallery or PDF reader. The PDF reader exposes page thumbnails and zoom for vouchers or instructions. The size experiment measures exact-byte deduplication in the beta.

\subsection{Viewer library and day navigation}
The Viewer opens the local trip library without a sign-in step. It imports a package from the file picker or a supported operating-system file-open action. A newly accepted itinerary appears in the local trip library. The traveler opens a trip, selects a day, and can swipe between days or select an event from the day index to scroll to that card. In the stable Viewer, selecting a day scrolls the timeline to its first card. Scrolling across a day boundary updates the selected day. The timeline labels times in each event's recorded zone and opens a card's location, photos, or documents in a dedicated view. Trip-wide material remains accessible from Information and Documents. The Viewer renders the schedule in read-only mode. Checking a packing item changes its local state against the trip and item identifiers. Schedule replacement preserves those records and the trip's completion status. The stable Viewer displays a container integrity indicator after envelope and asset checks complete. The beta supports active and completed views and lets the traveler restore a completed trip to the active list.

\subsection{Stable delivery options}
The stable Creator offers a printable PDF with the day schedule, booking references, and emergency contacts. It also offers editable templates alongside ordinary and password-protected exports. The template compiler removes traveler records, booking references, and private vouchers before packaging the remaining data. Reusing a template starts a new editable trip with its own identifier. The distribution portal publishes app downloads, release checksums, and a gallery of sanitized templates. Private trip delivery uses the agent's chosen transfer channel. Its roadmap also includes a browser Viewer, an anonymous release feed, tablet authoring, and direct local transfer. These capabilities define the wider product design. The experiments in this paper use the native Creator and Viewer code paths and the shared Dart core.
